\centerline{\bf \Huge Комплексные числа. }

\begin{flushright}
        \scriptsize{}
\end{flushright}

\textbf{Определение.} Комплексные числа --- это числа вида $z=x+i y$, где $x$ и $y$-действительные, а $i$ - мнимая единица, т.е. число, квадрат которого равен -1 . Числа $x, y$ называются соответственно действительной и мнимой частью. комплексного числа $z$ и обозначаются $x=\operatorname{Re} z, y=\operatorname{Im} z$. Сопряжённым числом к $z=x+i y$ называют $\bar{z}=x-i y$. Модулем числа $z=x+i y$ называют число $|z|=\sqrt{x^2+y^2}$. Множество всех комплексных чисел обозначается $\mathbb{C}$.

\textbf{Свойства сопряжения:}

(a) $\overline{z+w}=\bar{z}+\bar{w} ; \overline{z \cdot w}=\bar{z} \cdot \bar{w} ; \overline{z / w}=\bar{z} / \bar{w} ; \overline{\bar{z}}=z$;

(b) $|z|^2=z \bar{z} ; \operatorname{Re} z=\frac{z+\bar{z}}{2} ; \operatorname{Im} z=\frac{z-\bar{z}}{2 i}$.

\textbf{Тригонометрическая форма.}
Для любого комплексного числа $z \neq 0$ справедливо равенство $z=|z|(\cos \varphi+i \sin \varphi)$, где $|z|$ и $\varphi$ - соответственно модуль и аргумент числа $z$.

\textbf{Геометрический смысл умножения на комплексное число.}
При умножении (делении) двух комплексных чисел их аргументы складываются (вычитаются), а модули перемножаются (делятся). Отсюда нетрудно вывести формулу Муавра:

$$
(r(\cos \varphi+i \sin \varphi))^n=r^n(\cos n \varphi+i \sin n \varphi) .
$$

\newpage

\hspace{3mm}

\textbf{\Large Алгебраическая разминка}

\problem{1}
Используя формулу Муавра, вычислите

a) $(1+i)^n$;

б) $(1+i \sqrt{3})^n$;

в) $\left(\frac{1+i \sqrt{3}}{1-i}\right)^{20}$;

г) $\left(1-\frac{\sqrt{3}-i}{2}\right)^{20}$; д) $(1+\cos \varphi+i \sin\varphi)^n$;

e) $(\sqrt{3}+i)^n$;

ж) $\left(\dfrac{\cos \varphi+i \sin \varphi}{\cos \psi+i \sin \psi}\right)^n$.

\problem{2}
Решите уравнения:

(a) $z^2-5 z+4+10 i=0$;

(b) $z^3-1=0$.

\problem{3}
Найдите все комплексные числа сопряженные своему

(a) квадрату;

(b) кубу.

\problem{4}
Докажите равенство $|z+w|^2+|z-w|^2=2|z|^2+2|w|^2$.

\problem{5}
Вычислите суммы:

(a) $\sin \varphi+\cdots+\sin n \varphi$;

(b) $C_n^0-C_n^2+C_n^4-\ldots$;

(c) $C_n^0+C_n^4+C_n^8+\ldots$.


\newpage

\hspace{3mm}

\textbf{\Large Геометрическая тренировка}

\textbf{Деление отрезка в данном отношении.}
Если точка $C$ лежит на прямой $A B$ и $\overrightarrow{A C}=\lambda \overrightarrow{C B}(\lambda \neq-1)$, то говорят, что точка $C$ делит отрезок $A B$ в отношении $\lambda(\lambda=\bar{\lambda})$. Комплексные координаты точки $c$ записываются так:


$$
c=\frac{a+\lambda b}{1+\lambda}, \quad \lambda=\bar{\lambda} .
$$

\problem{1}
Даны два параллелограмма $A B C D$ и $A_1 B_1 C_1 D_1$. Докажите, что точки, делящие отрезки $A A_1, B B_1, C C_1, D D_1$ в одном и том же отношении, служат вершинами параллелограмма.

\problem{2}
Дан правильный треугольник $A B C$ и комплексная координата $a$ вершины $A$. Найдите комплексную координату вершины $B$ при положительной и отрицательной ориентациях треугольника $A B C$, если за начальную точку принята 1) вершина $C$, 2) центр треугольника $A B C$, 3) основание $A_1$ высоты $A A_1$.

\textbf{Расстояние между двумя точками} $A(a)$ и $B(b)$ вычисляется следующим образом:

$$
A B^2=(a-b)(\bar{a}-\bar{b}).
$$

\problem{3}
Средняя линия четырёхугольника делит его на два четырёхугольника. Докажите, что середины диагоналей этих двух четырёхугольников являются вершинами параллелограмма, либо лежат на одной прямой

\problem{4}
Докажите, что сумма квадратов диагоналей параллелограмма равна сумме квадратов всех его сторон.

\problem{5}
Докажите, что сумма квадратов диагоналей четырёхугольника равна удвоенной сумме квадратов его средних линий.

\newpage

\hspace{3mm}

\problem{6*} Дан треугольник $ABC$. Пусть точки касания вписанной окружности имеют комплексные координаты $p,q,r$. Известно, что точка Фейербаха имеет координату $\dfrac{pq+qr+pr}{p+q+r}$, а середины сторон имеют координаты $M_a = \dfrac{ab}{a+b}+\dfrac{ac}{a+c}, M_b = \dfrac{ab}{a+b} + \dfrac{bc}{b+c}, M_{c} = \dfrac{ac}{a+c} + \dfrac{bc}{b+c}$. Докажите, что из трёх отрезков $FM_{a}, FM_{b}, FM_{c}$ больший отрезок равен сумме двух других.